\chapter{Kết luận và hướng phát triển}
\label{chap:conclusion}

\section{Kết luận}

Khóa luận nghiên cứu bài toán học thói quen sinh hoạt của \emph{từng người} trong nhà thông minh có
nhiều người, chỉ từ cảm biến môi trường --- nơi các sự kiện của hai người trộn vào một luồng vô danh.
Xuất phát từ đề tài thực tập (tìm thói quen cho nhà một người) và định hướng từ luận án của
Chen~\cite{chen2021thesis} (nhận dạng hoạt động trong nhà nhiều người), khóa luận đi tiếp theo hướng
luận án còn bỏ ngỏ: từ nhận dạng hoạt động hiện tại sang hành vi dài hạn.

Các kết quả chính:

\begin{enumerate}
  \item \textbf{Phương pháp theo vùng của đề tài thực tập không dùng được cho nhà hai người.} Trên Cairo,
  hai người đi lại xen kẽ làm phiên bị cắt vụn, và DBSCAN gộp toàn bộ các bữa ăn cùng việc uống thuốc
  trong vùng bếp thành một cụm duy nhất. Tách luồng theo vị trí (B1) không khắc phục được, thậm chí còn
  làm kết quả tệ hơn.
  \item \textbf{Học (người, hoạt động) rồi áp định nghĩa thói quen cho kết quả tốt hơn hẳn.} Ở ba nhà hai
  người, F1 thói quen của Random Forest đạt 0,50--0,71 so với 0,22--0,33 của phương pháp thực tập; ở Cairo, Random Forest
  tìm đủ 6/6 thói quen gắn người, trong đó có thói quen uống thuốc của R2 --- đúng trường hợp mà cách
  tiếp cận theo vùng không thể tách được.
  \item \textbf{Định nghĩa thói quen được hiện thực đúng nguyên văn} bằng thuật toán quét cửa sổ đếm theo
  ngày. Ablation cho thấy DBSCAN và HDBSCAN lệch khỏi định nghĩa đáng kể ngay cả trên nhãn thật.
  \item \textbf{Thói quen chịu được lệch giờ nhỏ và nhận thừa ngẫu nhiên, nhưng nhạy với bỏ sót, nhầm
  người và lỗi có hệ thống.} Lỗi chính của các mô hình hiện tại là thói quen giả sinh từ lỗi nhận dạng lặp
  lại đều đặn, không phải bỏ sót.
  \item \textbf{Phát hiện ngày bỏ thuốc khả thi một phần}: Bi-GRU phát hiện 9/10 ngày bỏ thuốc được cài
  nhưng còn nhiều báo động nhầm; cách tiếp cận theo vùng không phát hiện được ngày nào. Mô hình tốt nhất cho
  thói quen (Random Forest) lại không phát hiện được sự vắng mặt --- một bài học về sự khác nhau giữa
  ``nhận ra điều đang có'' và ``nhận ra điều đã không xảy ra''.
\end{enumerate}

\section{Hạn chế}

\begin{itemize}
  \item \textbf{Dữ liệu}: chỉ có ba nhà hai người, nhãn gắn người ít (đánh giá ``từng người'' chủ yếu dựa
  trên Cairo), nhãn thưa và dữ liệu thu từ 2009--2011 ở Mỹ.
  \item \textbf{Đánh giá}: số ngày kiểm tra mỗi lần chia ít nên kết quả dao động lớn; đánh giá qua nhiều lần
  chia mới làm cho ba nhà hai người; mẫu của RQ3 nhỏ và bất thường là bất thường được cài.
  \item \textbf{Phương pháp}: bộ nhận dạng chưa được tối ưu cho mục tiêu thói quen (ít bỏ sót, ít lỗi có hệ
  thống); bước giải mã còn đơn giản.
\end{itemize}

\section{Hướng phát triển}

\textbf{Hướng phát triển trực tiếp của khóa luận:}
\begin{enumerate}
  \item Mở rộng đánh giá qua nhiều lần chia cho cả năm nhà và với dữ liệu dài hơn, để thu hẹp độ dao động.
  \item Giảm thói quen giả: giải mã có xét thời lượng hoạt động, ràng buộc ``mỗi người mỗi ngày tối đa một
  lần thức dậy'', hoặc hiệu chỉnh ngưỡng theo mục tiêu thói quen thay vì F1 theo sự kiện.
  \item Giảm phụ thuộc vào giờ trong ngày của bộ nhận dạng (Hình~\ref{fig:rf-imp}), ví dụ bằng đặc trưng
  mô tả chuỗi chuyển động thay cho giờ, để mô hình nhận ra cả sự \emph{vắng mặt} của một hoạt động.
  \item Mở rộng dữ liệu nếu xin được: ARAS~\cite{alemdar2013aras}, CASAS Kyoto/twor.
\end{enumerate}

\textbf{Ngoài phạm vi khóa luận} (các ứng dụng xây trên nền thói quen):
\begin{itemize}
  \item Hệ thống nhắc uống thuốc thật, có giao diện cho người chăm sóc.
  \item Dự đoán hoạt động tiếp theo dựa trên thói quen~\cite{doctorarastoo2026prediction}.
  \item Dùng mô hình ngôn ngữ để giải thích thói quen và thay đổi hành vi cho gia đình hoặc bác
  sĩ~\cite{civitarese2026xbcd}, hoặc mô tả cảm biến bằng văn bản để chuyển sang nhà
  mới~\cite{thukral2024tdost,fiori2026domusfm}.
  \item Nhà từ ba người trở lên với dữ liệu như MuRAL~\cite{chen2025mural}.
\end{itemize}
